\documentclass[11pt]{article}
\usepackage[margin=1in]{geometry}
\usepackage{listings,xcolor,hyperref,booktabs,longtable,amsmath,amssymb}
\lstset{basicstyle=\ttfamily\small,breaklines=true,frame=single,numbers=left,numberstyle=\tiny\color{gray},
language=C,morekeywords={kernel,__global__,__shared__,__device__}}
\title{Private LLM Kernel Acceleration \& Quality Controls:\\Zig--CUDA Hybrid Inference Engine}
\author{SAP OSS Technical Documentation}
\date{March 2026}

\begin{document}
\maketitle

\begin{abstract}
This report documents the architecture, kernel design, and quality controls of a
private LLM inference engine built as a pure Zig + CUDA pipeline. The system is
\emph{not} a fork of vLLM or any Python framework; it is a ground-up
implementation targeting NVIDIA T4 (sm\_75) and L4 (sm\_89) GPUs via Zig 0.15.1
with hand-written CUDA kernels compiled to PTX. The engine supports DeltaNet
hybrid architectures (Qwen3.5), speculative decoding with multi-token prediction,
and multi-backend GPU dispatch across CUDA, Metal, and WebGPU. We present
exact code evidence with file paths and line numbers throughout.
\end{abstract}

\tableofcontents
\newpage

% ============================================================================
\section{System Architecture}
\label{sec:architecture}
% ============================================================================

The inference engine is a pure Zig application with native CUDA kernel
integration. It does \textbf{not} depend on Python, PyTorch, or the vLLM
serving framework. The entire forward pass, memory management, and kernel
dispatch are implemented in Zig 0.15.1.

\subsection{Build System}

The top-level build configuration resides in:
\begin{center}
\texttt{src/intelligence/vllm-main/zig/build.zig}
\end{center}

The GPU build is invoked with architecture-specific flags:

\begin{lstlisting}[language=bash,caption={Build commands for T4 and L4 targets}]
# Tesla T4 (sm_75, 15GB VRAM, CUDA 13.0)
zig build -Doptimize=ReleaseFast -Dgpu=true -Dgpu_arch=sm_75

# NVIDIA L4 (sm_89, 23GB VRAM, CUDA 12.8)
zig build -Doptimize=ReleaseFast -Dgpu=true -Dgpu_arch=sm_89
\end{lstlisting}

The architecture flag is resolved at \texttt{build.zig} line 12:

\begin{lstlisting}[language=C,caption={GPU architecture option (build.zig:12)}]
// GPU architecture: sm_75=T4, sm_89=L4, sm_90=H200/H100.
// Defaults to sm_75.
const gpu_arch = b.option([]const u8, "gpu_arch",
    "CUDA GPU architecture (sm_75/sm_89/sm_90)")
    orelse "sm_75";
\end{lstlisting}

When \texttt{-Dgpu=true} is set, the build system invokes \texttt{nvcc} to
compile 16 CUDA source files into a shared library
(\texttt{build.zig} lines 19--48):

\begin{lstlisting}[language=C,caption={CUDA kernel compilation (build.zig:19--48)}]
cuda_kernels_build = b.addSystemCommand(&.{
    "nvcc", "-O3", arch_flag, "-I", "csrc",
    "-lcublas", "-shared", "-Xcompiler", "-fPIC",
    "csrc/cuda_kernels.cu",
    "csrc/continuous_batching.cu",
    "csrc/cuda_graphs.cu",
    "csrc/flash_attention.cu",
    "csrc/flash_attention_v2.cu",
    "csrc/fp8_quantization.cu",
    "csrc/fused_kernels.cu",
    "csrc/glm5_kernels.cu",
    "csrc/gpu_tokenizer.cu",
    "csrc/int8_kv_cache.cu",
    "csrc/int8_quantization.cu",
    "csrc/kimi25_kernels.cu",
    "csrc/minimax25_kernels.cu",
    "csrc/mla_kernels.cu",
    "csrc/pipeline_parallel.cu",
    "csrc/tensor_core_ops.cu",
    "csrc/tensor_parallel.cu",
    "-o", "zig-out/lib/libcuda_kernels.so",
});
\end{lstlisting}

\subsection{GPU Source File Inventory}

The \texttt{src/gpu/} directory contains 15 Zig modules plus 3 kernel sub-modules,
totaling 18 GPU source files:

\begin{longtable}{lp{3.2in}}
\toprule
\textbf{File} & \textbf{Purpose} \\
\midrule
\texttt{cuda\_forward.zig} & Full transformer forward pass on GPU \\
\texttt{cuda\_weights.zig} & GPU-resident quantized weight storage (Q4\_0, Q4\_K, Q5\_K) \\
\texttt{cuda\_backend.zig} & CUDA backend with PTX kernel discovery \\
\texttt{cuda\_bindings.zig} & CUDA driver API bindings \\
\texttt{cuda\_kernels.zig} & Zig-compiled PTX kernels \\
\texttt{backend.zig} & Unified GPU backend abstraction \\
\texttt{metal\_backend.zig} & Apple Metal compute backend \\
\texttt{metal\_shaders.zig} & Metal shader library management \\
\texttt{metal\_bindings.zig} & Metal Objective-C runtime bindings \\
\texttt{webgpu\_backend.zig} & WebGPU/Vulkan via wgpu-native \\
\texttt{accelerate\_backend.zig} & Apple Accelerate framework (BLAS/vDSP) \\
\texttt{memory\_pool.zig} & Triple-buffered GPU memory allocator \\
\texttt{zero\_copy\_pipeline.zig} & Zero-copy HTTP-to-GPU data transfer \\
\texttt{async\_pipeline.zig} & Asynchronous compute/transfer overlap \\
\texttt{context.zig} & GPU context and buffer management \\
\texttt{kernels/kernel\_fusion.zig} & Fused JSON$\to$tokenize$\to$inference pipeline \\
\texttt{kernels/json\_parser.zig} & GPU-side JSON field extraction \\
\texttt{kernels/tokenizer.zig} & GPU-native BPE/SentencePiece tokenizer \\
\bottomrule
\end{longtable}

Additionally, the CUDA-specific kernel files reside alongside:
\begin{itemize}
  \item \texttt{deltanet\_kernels.cu} --- 588 lines of CUDA C source
  \item \texttt{deltanet\_kernels.ptx} --- 2441 lines of compiled PTX assembly
\end{itemize}

\subsection{Module Dependency Graph}

The build system links the following module chain (\texttt{build.zig} lines 113--162):

\begin{enumerate}
  \item \texttt{cuda\_forward\_mod} imports \texttt{cuda\_build\_options} (line 118)
  \item \texttt{main\_mod} imports \texttt{llama}, \texttt{metal\_bindings}, \texttt{metal\_shaders}, and \texttt{cuda\_build\_options} (lines 149--162)
  \item CUDA modules use relative \texttt{@import} within \texttt{src/gpu/} to avoid file-ownership conflicts (comment at lines 158--161)
  \item macOS builds link Metal, Foundation, CoreGraphics, and Accelerate frameworks (lines 177--182)
\end{enumerate}

% ============================================================================
\section{DeltaNet Hybrid Architecture}
\label{sec:deltanet}
% ============================================================================

The engine supports the Qwen3.5 hybrid architecture which interleaves
standard multi-head attention layers with DeltaNet linear attention layers.
Configuration and state management are in:
\begin{center}
\texttt{src/intelligence/vllm-main/zig/src/gpu/cuda\_forward.zig} (lines 149--416)
\end{center}

\subsection{CudaForwardConfig Hybrid Fields}

The \texttt{CudaForwardConfig} struct (line 124) extends the standard transformer
configuration with DeltaNet-specific fields at lines 149--157:

\begin{lstlisting}[language=C,caption={DeltaNet hybrid fields (cuda\_forward.zig:149--157)}]
// Gated DeltaNet / hybrid fields (Qwen3.5)
// zero = pure transformer
ssm_inner_size: u32 = 0,     // DeltaNet V heads * state_size (e.g. 2048)
ssm_state_size: u32 = 0,     // DeltaNet head_dim (e.g. 128)
ssm_group_count: u32 = 0,    // DeltaNet QK heads (e.g. 16)
ssm_conv_kernel: u32 = 0,    // conv1d kernel size (e.g. 4)
ssm_time_step_rank: u32 = 0, // num_heads for decay (e.g. 16)
attn_head_dim: u32 = 0,      // attention head_dim (e.g. 256)
rope_dim: u32 = 0,           // partial RoPE dimension (e.g. 64)
full_attn_interval: u32 = 0, // every Nth layer is full attention
\end{lstlisting}

\subsection{Layer Type Dispatch Methods}

Three key methods determine the execution path per layer (lines 164--177):

\begin{lstlisting}[language=C,caption={Hybrid dispatch methods (cuda\_forward.zig:164--177)}]
/// True if hybrid DeltaNet+Attention architecture (Qwen3.5)
pub fn isHybrid(self: CudaForwardConfig) bool {
    return self.full_attn_interval > 0;
}

/// True if layer `l` is a full attention layer (vs DeltaNet)
pub fn isAttnLayer(self: CudaForwardConfig, l: u32) bool {
    if (self.full_attn_interval == 0) return true; // pure transformer
    return (l + 1) % self.full_attn_interval == 0;
}

/// Number of DeltaNet V/output heads (= time_step_rank)
pub fn ssmNumHeads(self: CudaForwardConfig) u32 {
    return self.ssm_time_step_rank;
}
\end{lstlisting}

The \texttt{isAttnLayer} method implements the interleaved pattern: given
\texttt{full\_attn\_interval = 4}, layers 3, 7, 11, \ldots use standard
multi-head attention, while all others use DeltaNet linear attention.

\subsection{Dimension Computation Helpers}

Lines 179--217 provide dimension accessors for the asymmetric Q/K/V structure:

\begin{itemize}
  \item \texttt{ssmVDim()} = \texttt{time\_step\_rank} $\times$ \texttt{state\_size} (V dimension, line 181)
  \item \texttt{ssmKVDim()} = \texttt{group\_count} $\times$ \texttt{state\_size} (KV dimension, line 186)
  \item \texttt{ssmQKVDim()} = $2 \times$ \texttt{ssmKVDim()} + \texttt{ssmVDim()} (total projection, line 191)
  \item \texttt{attnHeadDim()} returns \texttt{attn\_head\_dim} or fallback (line 201)
  \item \texttt{ropeDim()} returns partial RoPE dimension or full head\_dim (line 206)
\end{itemize}

\subsection{DeltaNetState: Persistent Recurrent State}

The \texttt{DeltaNetState} struct (lines 296--416) manages GPU-resident recurrent
state that persists across tokens within a sequence:

\begin{lstlisting}[language=C,caption={DeltaNetState layout (cuda\_forward.zig:296--316)}]
const DeltaNetState = struct {
    n_dn_layers: u32,   // number of DeltaNet layers
    num_v_heads: u32,   // Q/output heads (= time_step_rank)
    num_kv_heads: u32,  // KV heads (= group_count)
    state_size: u32,    // head_dim for DeltaNet (e.g. 128)
    conv_kernel: u32,
    ssm_inner: u32,     // ssm_inner_size (V heads * state_size)
    kv_dim: u32,        // group_count * state_size
    qkv_dim: u32,       // ssm_inner + 2 * kv_dim
    num_heads: u32,     // time_step_rank

    // Persistent state (survives across tokens)
    d_S: cuda.CUdeviceptr = 0,    // [n_dn * V_heads * D * D] f32
    d_conv: cuda.CUdeviceptr = 0, // [n_dn * (K-1) * qkv_dim] f32

    // Scratch buffers (reused each token)
    d_qkv: cuda.CUdeviceptr = 0,   // [qkv_dim] f32
    d_gate: cuda.CUdeviceptr = 0,  // [ssm_inner] f32
    d_alpha: cuda.CUdeviceptr = 0, // [num_heads] f32
    d_beta: cuda.CUdeviceptr = 0,  // [num_heads] f32
    d_y: cuda.CUdeviceptr = 0,     // [ssm_inner] f32
};
\end{lstlisting}

\subsubsection{GPU Memory Layout}

The per-layer S matrix has layout:
\[
S[\text{layer}][\text{head}][d_1][d_2] \in \mathbb{R}^{n_{\text{dn}} \times n_v \times D \times D}
\]
where $D = \texttt{state\_size}$. Per-layer byte computation (line 318):

\begin{lstlisting}[language=C,caption={Per-layer S matrix bytes (cuda\_forward.zig:318--320)}]
fn sLayerBytes(self: DeltaNetState) usize {
    return @as(usize, self.num_v_heads)
         * self.state_size * self.state_size * @sizeOf(f32);
}
\end{lstlisting}

Layer-indexed device pointers (lines 327--334):

\begin{lstlisting}[language=C,caption={Device pointer indexing (cuda\_forward.zig:327--334)}]
fn sPtr(self: DeltaNetState, dn_layer: usize) cuda.CUdeviceptr {
    return self.d_S + dn_layer * self.sLayerBytes();
}

pub fn convPtr(self: DeltaNetState, dn_layer: usize) cuda.CUdeviceptr {
    return self.d_conv + dn_layer * self.convLayerBytes();
}
\end{lstlisting}

\subsubsection{State Lifecycle}

\begin{itemize}
  \item \textbf{Init} (line 336): Counts DeltaNet layers, allocates S and conv
        state via \texttt{cuMemAlloc}, zero-initializes with \texttt{cuMemsetD32}.
  \item \textbf{Clear} (line 410): Resets persistent state to zero for new sequences
        via \texttt{cuMemsetD32}.
  \item \textbf{Deinit} (line 396): Frees all GPU buffers via \texttt{cuMemFree},
        nulls pointers.
\end{itemize}

% ============================================================================
\section{CUDA Kernels}
\label{sec:kernels}
% ============================================================================

The DeltaNet-specific kernels are in two files:
\begin{itemize}
  \item \texttt{src/intelligence/vllm-main/zig/src/gpu/deltanet\_kernels.cu} (588 lines, CUDA C)
  \item \texttt{src/intelligence/vllm-main/zig/src/gpu/deltanet\_kernels.ptx} (2441 lines, PTX assembly)
\end{itemize}

The PTX was generated by NVVM 7.0.1, CUDA 12.1, targeting \texttt{sm\_75}
(PTX header, lines 1--11):

\begin{lstlisting}[language=C,caption={PTX header (deltanet\_kernels.ptx:1--11)}]
// Generated by NVIDIA NVVM Compiler
// Compiler Build ID: CL-32688072
// Cuda compilation tools, release 12.1, V12.1.105
// Based on NVVM 7.0.1
.version 8.1
.target sm_75
.address_size 64
\end{lstlisting}

\subsection{Kernel Entry Points}

The PTX contains 12 kernel entry points, discovered via \texttt{.visible .entry}
declarations:

\begin{longtable}{rlp{2.8in}}
\toprule
\textbf{PTX Line} & \textbf{Kernel} & \textbf{Function} \\
\midrule
18 & \texttt{q4\_gemv} & Q4\_0 quantized GEMV with shared memory \\
406 & \texttt{q4\_1\_gemv} & Q4\_1 quantized GEMV (min+scale) \\
801 & \texttt{deltanet\_conv1d} & Depthwise conv1d with state shift \\
1017 & \texttt{deltanet\_l2norm} & Per-head L2 normalization \\
1122 & \texttt{deltanet\_gates} & Alpha (decay) and beta (update) gates \\
1279 & \texttt{deltanet\_recurrent} & Core recurrent state update + readout \\
1583 & \texttt{deltanet\_output\_gate} & RMSNorm $\times$ SiLU gating \\
1725 & \texttt{partial\_rope\_q} & Partial RoPE on Q (64 of 256 dims) \\
1829 & \texttt{partial\_rope\_k} & Partial RoPE on K \\
1933 & \texttt{split\_q\_gate} & Split fused Q+gate for hybrid attention \\
1984 & \texttt{decode\_attention} & GQA decode attention (arbitrary head\_dim) \\
2381 & \texttt{gated\_attn\_output} & Gated sigmoid attention output \\
\bottomrule
\end{longtable}

\subsection{Q4\_0 Quantized GEMV}

The \texttt{q4\_gemv} kernel (\texttt{deltanet\_kernels.cu} lines 16--73)
implements $\mathbf{y}[M] = \mathbf{W}_{q4}[M \times K] \cdot \mathbf{x}[K]$
with the following optimizations:

\begin{enumerate}
  \item \textbf{Cooperative shared memory load}: Input vector $\mathbf{x}$ loaded via
        128-bit coalesced reads (\texttt{float4}) into shared memory (lines 26--33).
  \item \textbf{Hoisted bias}: The $-8$ offset is folded into a pre-multiplied
        \texttt{neg8s = -8.0f * scale}, eliminating per-element subtraction (line 51).
  \item \textbf{Warp shuffle reduction}: Butterfly reduction via
        \texttt{\_\_shfl\_xor\_sync} avoids shared memory bank conflicts (lines 68--69).
\end{enumerate}

\begin{lstlisting}[caption={Q4\_0 GEMV core loop (deltanet\_kernels.cu:46--69)}]
for (int b = lane; b < n_blocks; b += 32) {
    const uint8_t* block = W_row + b * 18;
    float scale = __half2float(
        __ldg(reinterpret_cast<const __half*>(block)));
    float neg8s  = -8.0f * scale;
    int x_base = b << 5;
    const uint8_t* data = block + 2;
    #pragma unroll
    for (int j = 0; j < 16; j++) {
        uint32_t bv = __ldg(&data[j]);
        acc += ((float)(int)(bv & 0xF) * scale + neg8s)
               * x_smem[x_base + j];
        acc += ((float)(int)(bv >> 4) * scale + neg8s)
               * x_smem[x_base + j + 16];
    }
}
// Warp shuffle butterfly reduction
for (int offset = 16; offset > 0; offset >>= 1)
    acc += __shfl_xor_sync(0xFFFFFFFF, acc, offset);
\end{lstlisting}

Grid configuration: $(\lceil M/8 \rceil, 1)$ blocks of $(256, 1)$ threads ---
8 warps per block, 1 warp per output row.

\subsection{DeltaNet Recurrent Kernel}

The core recurrent update (\texttt{deltanet\_kernels.cu} lines 283--341) implements
the gated delta rule with optimized memory access:

\[
S_{t} = \alpha \cdot S_{t-1} + \beta \cdot \mathbf{k} \otimes (\mathbf{v} - \mathbf{k}^\top S_{t-1})
\]
\[
\mathbf{y} = S_t^\top \mathbf{q}
\]

Key optimizations over the naive 3-pass approach:

\begin{enumerate}
  \item K and Q vectors loaded into shared memory once (lines 301--312),
        eliminating repeated global reads.
  \item Three logical passes (decay, write, readout) fused into two physical
        passes: decay+retrieve in pass 1, write+readout in pass 2
        (lines 321--340).
\end{enumerate}

\begin{lstlisting}[caption={DeltaNet recurrent core (deltanet\_kernels.cu:319--340)}]
for (int d2 = threadIdx.x; d2 < D; d2 += blockDim.x) {
    // Pass 1: decay S[d1][d2], accumulate kv_mem
    float kv_mem = 0.0f;
    for (int d1 = 0; d1 < D; d1++) {
        float s = S_h[d1 * D + d2] * a;
        S_h[d1 * D + d2] = s;
        kv_mem += s * k_smem[d1];
    }
    // Delta
    float delta = (v_h[d2] - kv_mem) * b;
    // Pass 2: write delta + readout (fused)
    float y_val = 0.0f;
    for (int d1 = 0; d1 < D; d1++) {
        float s = S_h[d1 * D + d2] + k_smem[d1] * delta;
        S_h[d1 * D + d2] = s;
        y_val += s * q_smem[d1];
    }
    y_h[d2] = y_val;
}
\end{lstlisting}

\subsection{DeltaNet Gate Computation}

The gate kernel (\texttt{deltanet\_kernels.cu} lines 244--267) computes:
\begin{align}
\alpha_h &= \exp\!\bigl(A_h \cdot \text{softplus}(\text{proj}_h + \text{bias}_h)\bigr) \\
\beta_h &= \sigma(\text{proj}_h)
\end{align}
where $A_h = -\exp(A_{\log,h})$ is pre-negated in the GGUF weight converter.

\subsection{Partial RoPE}

For Qwen3.5 attention layers with \texttt{head\_dim=256} and \texttt{rope\_dim=64},
only the first 64 elements per head receive rotary position embeddings
(\texttt{deltanet\_kernels.cu} lines 398--461). The kernel uses a shared-memory
cos/sin lookup table built via a \texttt{\_\_device\_\_} helper
(\texttt{build\_rope\_table}, lines 407--417).

\subsection{PTX Instruction Highlights}

The compiled PTX uses several performance-critical ISA features:
\begin{itemize}
  \item \texttt{ld.global.nc.v4.f32} --- Non-coherent texture cache loads for
        read-only data (PTX line 50).
  \item \texttt{fma.rn.f32} --- Fused multiply-add with round-to-nearest,
        used throughout Q4 dequantization.
  \item \texttt{bar.sync} --- Warp/block synchronization barriers for shared
        memory coordination.
  \item \texttt{shfl.sync.bfly} --- Butterfly shuffle for warp-level reductions
        without shared memory.
\end{itemize}

% ============================================================================
\section{Speculative Decoding}
\label{sec:spec-decode}
% ============================================================================

Speculative decoding is implemented in:
\begin{center}
\texttt{src/intelligence/vllm-main/zig/src/toon/llama\_toon.zig} (lines 809--882)
\end{center}

The system employs a dual-strategy draft generation pipeline with adaptive
disabling.

\subsection{Draft Generation Strategies}

Two complementary strategies are tried in order (lines 815--819):

\begin{lstlisting}[language=C,caption={Dual draft strategy (llama\_toon.zig:815--819)}]
var n_drafted = self.tryNativeMtpDraft(
    fwd, next_token, pos - 1, draft_buf[1..K]);
if (n_drafted == 0) {
    n_drafted = self.ngramDraft(
        context_buf.items, draft_buf[1..K]);
}
if (n_drafted == 0) break :spec_path;
\end{lstlisting}

\begin{enumerate}
  \item \textbf{Native MTP Draft} (\texttt{tryNativeMtpDraft}, lines 942--979):
        Uses the model's own multi-token prediction heads to generate draft tokens.
        Downloads hidden states from GPU (\texttt{fwd.activations.hidden.downloadF32}),
        then runs \texttt{drafter.fillContinuation} for up to
        \texttt{drafter.supportedPositions()} tokens.

  \item \textbf{N-gram Table Fallback} (\texttt{ngramDraft}, lines 665--686):
        Uses a hash-based trigram + bigram table to predict continuations.
        Tries trigram lookup first, falls back to bigram:
\begin{lstlisting}[language=C,caption={N-gram draft logic (llama\_toon.zig:674--678)}]
const pred = if (have_trigram or count > 0)
    (self.ngramLookup(t0, t1, t2)
     orelse self.bigramLookup(t1, t2))
else
    self.bigramLookup(t1, t2);
\end{lstlisting}
\end{enumerate}

\subsection{Batch Verification}

All $K$ draft tokens (including the verified token) are verified in a single
batch forward pass via \texttt{forwardBatchMoE} (lines 836--840):

\begin{lstlisting}[language=C,caption={Batch verification (llama\_toon.zig:836--840)}]
fwd.forwardBatchMoE(
    draft_buf[0..total_draft],
    positions_buf[0..total_draft],
    bl[0 .. @as(usize, total_draft) * vocab_size],
) catch break :spec_path;
\end{lstlisting}

\subsection{Acceptance Criterion}

The engine uses a strict \textbf{argmax acceptance} criterion (lines 844--852):
draft token $i$ is accepted if and only if \texttt{argmax(logits[i-1])} equals
\texttt{draft\_buf[i]}. Acceptance is sequential --- the first rejection
terminates the chain:

\begin{lstlisting}[language=C,caption={Argmax acceptance (llama\_toon.zig:844--852)}]
var num_accepted: u32 = 1;
for (1..total_draft) |i| {
    const predicted = argmaxLogits(
        bl[(i-1) * vocab_size .. i * vocab_size]);
    if (predicted == draft_buf[i]) {
        num_accepted += 1;
    } else {
        break;
    }
}
\end{lstlisting}

\subsection{Adaptive Disabling}

If no draft tokens are accepted after the first cycle, speculation is disabled
for the remainder of the request (lines 822--824):

\begin{lstlisting}[language=C,caption={Adaptive disable (llama\_toon.zig:822--824)}]
if (spec_cycles >= 1 and spec_accepted_total == 0) {
    break :spec_path;
}
\end{lstlisting}

\subsection{Position Tracking and Logits Reuse}

After acceptance, the position counter advances by \texttt{num\_accepted}
(line 869), and the logits from the last accepted token's batch output are
reused for the next sampling step (line 868):

\begin{lstlisting}[language=C,caption={Logits reuse (llama\_toon.zig:868--869)}]
logits = bl[(num_accepted-1) * vocab_size
            .. num_accepted * vocab_size];
pos += num_accepted;
\end{lstlisting}

% ============================================================================
\section{Inference Pipeline}
\label{sec:inference}
% ============================================================================

The full token generation loop resides in \texttt{llama\_toon.zig} method
\texttt{generateFromTokens} (lines 720--940).

\subsection{Prefill Phase}

Prompt tokens are fed sequentially through the GPU forward pass (lines 735--744):

\begin{lstlisting}[language=C,caption={GPU prefill (llama\_toon.zig:735--744)}]
if (use_gpu) {
    const fwd = self.cuda_forward.?;
    fwd.reset();
    for (tokens, 0..) |tok, p| {
        logits = fwd.forward(tok, p) catch |err| {
            std.log.err("GPU forward failed at pos={}", .{p});
            return err;
        };
    }
}
\end{lstlisting}

\subsection{Decode Loop}

The main decode loop (lines 788--913) supports both GPU and CPU paths:

\begin{lstlisting}[language=C,caption={GPU/CPU forward dispatch (llama\_toon.zig:886--892)}]
if (use_gpu) {
    logits = try fwd.forward(next_token, pos);
    self.native_mtp_hidden_valid =
        self.native_mtp_drafter != null;
} else {
    logits = self.model.?.forward(
        next_token, pos, &self.kv_cache.?);
}
pos += 1;
\end{lstlisting}

\subsection{EOS Detection}

EOS checking supports both GGUF-embedded and custom tokenizers (line 803):

\begin{lstlisting}[language=C,caption={Dual EOS check (llama\_toon.zig:803)}]
const is_eos = if (self.gguf_tokenizer) |gt|
    gt.isEos(next_token)
else
    self.tokenizer.?.isEos(next_token);
if (is_eos) break;
\end{lstlisting}

EOS is also checked within accepted speculative draft tokens (lines 861--862).

\subsection{Performance Timing}

The engine tracks three timing metrics:
\begin{itemize}
  \item \texttt{first\_token\_step\_ns} --- Time to first output token (line 780)
  \item \texttt{decode\_step\_sum\_ns} --- Cumulative decode step time (line 779)
  \item \texttt{avg\_step\_us} --- Per-token average (line 919)
\end{itemize}

For speculative decoding, additional metrics are logged (lines 927--933):
acceptance rate (\texttt{alpha\_pct}), effective TPS (\texttt{eff\_tps}),
total drafted vs.\ accepted tokens.

% ============================================================================
\section{GPU Memory Subsystem}
\label{sec:memory}
% ============================================================================

\subsection{Memory Pool (memory\_pool.zig)}

File: \texttt{src/intelligence/vllm-main/zig/src/gpu/memory\_pool.zig}

The \texttt{PoolConfig} struct (line 14) configures triple-buffered GPU memory:

\begin{lstlisting}[language=C,caption={Memory pool config (memory\_pool.zig:14--25)}]
pub const PoolConfig = struct {
    num_slots: usize = 3,              // triple buffering
    slot_size_bytes: usize = 64*1024*1024, // 64MB per slot
    max_elements_per_slot: usize = 1024,
    element_stride: usize = 4096,      // 1024 floats * 4B
    use_pinned_memory: bool = true,
};
\end{lstlisting}

Each \texttt{BufferSlot} (line 31) tracks its lifecycle through 5 states:
\texttt{free} $\to$ \texttt{cpu\_filling} $\to$ \texttt{ready\_for\_transfer}
$\to$ \texttt{gpu\_processing} $\to$ \texttt{gpu\_complete}.

\subsection{GPU Weight Storage (cuda\_weights.zig)}

File: \texttt{src/intelligence/vllm-main/zig/src/gpu/cuda\_weights.zig}

The \texttt{GGMLType} enum (line 26) supports 13 quantization formats from F32
down to Q2\_K. The Q4\_0 format uses 18 bytes per 32 values (line 56):

\begin{lstlisting}[language=C,caption={Q4\_0 block format (cuda\_weights.zig:5--9)}]
// Q4_0 block format (18 bytes per 32 values):
//   f16 d;       // scale factor
//   u8 qs[16];   // 32 nibbles packed into 16 bytes
//   dequant: val = (nibble - 8) * d
\end{lstlisting}

For a 7B parameter model in Q4\_0: $\approx$3.6 GB VRAM (vs.\ $\approx$28 GB in FP32).

\subsection{Zero-Copy Pipeline (zero\_copy\_pipeline.zig)}

File: \texttt{src/intelligence/vllm-main/zig/src/gpu/zero\_copy\_pipeline.zig}

Data flow bypasses CPU entirely:
\[
\text{HTTP Body} \xrightarrow{\text{DMA}} \text{VRAM}
\xrightarrow{\text{GPU Parse}} \text{GPU Tokenize}
\xrightarrow{} \text{GPU Inference} \to \text{Response}
\]

The \texttt{ZeroCopySlot} (line 41) uses 4096-byte aligned pinned memory for
DMA transfers and tracks 6 lifecycle states from \texttt{idle} through
\texttt{complete} (lines 69--76).

\subsection{Async Pipeline (async\_pipeline.zig)}

File: \texttt{src/intelligence/vllm-main/zig/src/gpu/async\_pipeline.zig}

Overlaps four stages using triple-buffered command slots:
\begin{enumerate}
  \item CPU batch preparation (\texttt{cpu\_filling})
  \item Host-to-device transfer (\texttt{h2d\_transfer})
  \item GPU compute (\texttt{gpu\_compute})
  \item Device-to-host transfer (\texttt{d2h\_transfer})
\end{enumerate}

Each \texttt{CommandSlot} (line 47) contains hardware-mapped \texttt{GpuBuffer}
instances for input tokens and output embeddings, with atomic state tracking
and \texttt{ResetEvent} completion signaling.

% ============================================================================
\section{Multi-Backend Support}
\label{sec:backends}
% ============================================================================

The engine abstracts GPU compute behind a unified interface supporting four
backends. All backend implementations reside in
\texttt{src/intelligence/vllm-main/zig/src/gpu/}.

\subsection{Backend Abstraction (backend.zig)}

The \texttt{GpuBackend} tagged union (line 49) dispatches to the active backend:

\begin{lstlisting}[language=C,caption={Unified backend dispatch (backend.zig:49--66)}]
pub const GpuBackend = union(BackendType) {
    metal: *MetalBackend,
    webgpu: *WebGPUBackend,
    cuda: *CudaBackend,
    cpu: *CpuFallback,

    pub fn submitBatch(self: GpuBackend,
                       batch: *const Batch) !BatchResult {
        return switch (self) {
            .metal  => |m| try m.submitBatch(batch),
            .webgpu => |w| try w.submitBatch(batch),
            .cuda   => |c| try c.submitBatch(batch),
            .cpu    => |f| try f.submitBatch(batch),
        };
    }
};
\end{lstlisting}

\texttt{BackendCapabilities} (line 33) reports hardware features including
\texttt{supports\_fp16}, \texttt{supports\_int8}, \texttt{supports\_async\_compute},
and \texttt{unified\_memory}.

\subsection{CUDA Backend (cuda\_backend.zig)}

Targets NVIDIA T4 (SM75), A100 (SM80), H100 (SM90). Features compile-time PTX
kernel discovery (lines 52--77): the embedded PTX binary is scanned at
\texttt{comptime} for \texttt{.visible .entry} declarations, populating a
kernel registry without runtime parsing.

\begin{lstlisting}[language=C,caption={Comptime PTX discovery (cuda\_backend.zig:52--77)}]
fn discoverPtxKernels() PtxKernelDiscovery {
    const ptx = @embedFile("cuda_kernels_ptx.bin");
    const marker = ".visible .entry ";
    var result = PtxKernelDiscovery{ .names = undefined,
                                     .count = 0 };
    @setEvalBranchQuota(500_000);
    while (pos + marker.len < ptx.len) : (pos += 1) {
        if (std.mem.startsWith(u8, ptx[pos..], marker)) {
            // extract kernel name...
            result.names[result.count] = ptx[name_start..name_end];
            result.count += 1;
        }
    }
    return result;
}
\end{lstlisting}

\subsection{Metal Backend (metal\_backend.zig)}

Targets Apple Silicon (M1/M2/M3) and Intel+AMD GPUs on macOS. Uses
triple-buffered command submission with Metal Performance Shaders (MPS)
support. Shared memory mode for Apple Silicon provides unified CPU/GPU
addressing. Configuration (line 34): up to 3 in-flight buffers at 64 MB each.

\subsection{WebGPU Backend (webgpu\_backend.zig)}

Cross-platform backend using wgpu-native (Vulkan on Linux/AWS). Enabled via
\texttt{zig build -Dwebgpu=true}. Extern declarations bind to 28 wgpu-native
C functions (lines 26--53) for instance creation, buffer management, compute
pipeline dispatch, and synchronization.

\subsection{Apple Accelerate Backend (accelerate\_backend.zig)}

Provides SIMD-optimized BLAS via the macOS Accelerate framework. Key functions:
\begin{itemize}
  \item \texttt{cblas\_sgemm} --- Single-precision matrix multiply (line 17)
  \item \texttt{cblas\_sgemv} --- Matrix-vector multiply (line 35)
  \item \texttt{vDSP\_vadd} --- Element-wise vector add (line 52)
\end{itemize}

\subsection{Metal Shader Infrastructure}

\begin{itemize}
  \item \texttt{metal\_shaders.zig} --- Shader library management and pipeline
        state caching
  \item \texttt{metal\_bindings.zig} --- Objective-C runtime bindings for Metal
        C API types (\texttt{MTLDevice}, \texttt{MTLCommandQueue}, etc.)
\end{itemize}

Build-time framework linking (\texttt{build.zig} lines 177--182):
\begin{lstlisting}[language=C,caption={macOS framework linking (build.zig:177--182)}]
if (target.result.os.tag == .macos) {
    exe.linkFramework("Metal");
    exe.linkFramework("Foundation");
    exe.linkFramework("CoreGraphics");
    exe.linkFramework("Accelerate");
}
\end{lstlisting}

% ============================================================================
\section{Kernel Fusion}
\label{sec:fusion}
% ============================================================================

The kernel fusion engine eliminates CPU round-trips by chaining three GPU stages
in a single dispatch. All files are in:
\begin{center}
\texttt{src/intelligence/vllm-main/zig/src/gpu/kernels/}
\end{center}

\subsection{Fusion Pipeline (kernel\_fusion.zig)}

The three-stage pipeline (\texttt{kernel\_fusion.zig} lines 1--9):

\begin{lstlisting}[language=C,caption={Kernel fusion architecture (kernel\_fusion.zig:4--9)}]
//! Stage 1: JSON Parser kernel extracts prompt text bounds
//! Stage 2: Tokenizer kernel converts text to tokens
//!          (no CPU roundtrip)
//! Stage 3: Inference kernel generates output
//!          (no CPU roundtrip)
//! All stages execute in GPU memory without returning to CPU
\end{lstlisting}

The \texttt{FusedResult} struct (line 47) captures per-stage timing and error
information:

\begin{lstlisting}[language=C,caption={Fused result structure (kernel\_fusion.zig:47--70)}]
pub const FusedResult = extern struct {
    json_status: u32,    text_start: u32,
    text_end: u32,       json_time_ns: u64,
    token_status: u32,   num_tokens: u32,
    token_time_ns: u64,
    inference_status: u32, output_tokens: u32,
    inference_time_ns: u64,
    total_time_ns: u64,
    error_stage: u32,    // 0=none, 1=parser,
                         // 2=tokenizer, 3=inference
    error_code: u32,
};
\end{lstlisting}

\subsection{GPU JSON Parser (json\_parser.zig)}

Pattern-based field extraction for OpenAI-compatible requests. The
\texttt{GpuParseResult} (line 24) returns byte offsets (\texttt{text\_start},
\texttt{text\_end}) directly readable by subsequent GPU kernels.
\texttt{MultiMessageParseResult} (line 40) supports up to 64 messages with
per-message role and content bounds.

\subsection{GPU Tokenizer (tokenizer.zig)}

Zero-CPU-touch tokenization supporting four algorithms:
\begin{itemize}
  \item BPE (GPT-2/GPT-3 style)
  \item WordPiece (BERT style)
  \item SentencePiece (T5/LLaMA style)
  \item Tiktoken (OpenAI style)
\end{itemize}

The tokenizer receives text bounds from the JSON parser and outputs token IDs
directly to GPU buffer, triggering inference via kernel fusion.

% ============================================================================
\section{Mangle A2A Rules}
\label{sec:mangle}
% ============================================================================

The agent-to-agent reasoning layer uses Mangle (a Datalog variant) for
declarative service routing and model selection. Rule files are in:
\begin{center}
\texttt{src/intelligence/vllm-main/mangle/a2a/}
\end{center}

\subsection{Fact Declarations (facts.mg)}

The fact schema (\texttt{facts.mg}, 141 lines) declares predicates without
hardcoded ground facts (loaded from config at runtime):

\begin{itemize}
  \item \texttt{service\_registry(name, base\_url, default\_model)} --- Service mesh entries
  \item \texttt{service\_status(name, status)} --- Temporal availability atoms (\texttt{/available}, \texttt{/unavailable}, \texttt{/degraded})
  \item \texttt{model\_registry(model\_name, service\_name, context\_length)} --- Model catalog
  \item \texttt{model\_specialization(model\_name, specialization)} --- Capability tags (\texttt{/code}, \texttt{/analysis}, \texttt{/reasoning})
  \item \texttt{agui\_event\_type(event\_type)} --- AG-UI protocol events (11 event types, lines 128--140)
\end{itemize}

\subsection{Derived Rules (rules.mg)}

The rules file (\texttt{rules.mg}, 212 lines) defines derived predicates:

\subsubsection{Service Routing}

\begin{lstlisting}[language=Prolog,caption={Intent resolution (rules.mg:25--29)}]
Decl resolve_service_for_intent(intent_id, url) :-
  intent_service_mapping(intent_id, service_name, endpoint_path),
  routable(service_name),
  service_registry(service_name, base_url, _),
  url = fn:concat(base_url, endpoint_path).
\end{lstlisting}

With fallback routing when the primary service is unavailable (lines 32--38):

\begin{lstlisting}[language=Prolog,caption={Fallback routing (rules.mg:32--38)}]
Decl resolve_fallback_service(intent_id, url) :-
  intent_service_mapping(intent_id, primary_service, _),
  !routable(primary_service),
  service_capability(fallback_service, /chat),
  routable(fallback_service),
  service_registry(fallback_service, base_url, _),
  url = fn:concat(base_url, "/v1/chat/completions").
\end{lstlisting}

\subsubsection{Model Selection}

\begin{lstlisting}[language=Prolog,caption={Model selection rules (rules.mg:109--118)}]
Decl best_model_for(specialization, model_name, service_name) :-
  model_specialization(model_name, specialization),
  model_registry(model_name, service_name, _),
  routable(service_name).

Decl model_for_context(min_context, model_name, service_name) :-
  model_registry(model_name, service_name, context_length),
  context_length >= min_context,
  routable(service_name).
\end{lstlisting}

\subsubsection{Quality Gates}

Service health monitoring via response latency (lines 91--102):

\begin{lstlisting}[language=Prolog,caption={Health quality gates (rules.mg:91--102)}]
Decl service_healthy(service_name) :-
  api_response(req_id, _, latency_ms, _),
  api_request(req_id, service_name, _, _, _),
  latency_ms < 1000.

Decl service_slow(service_name, avg_latency) :-
  service_registry(service_name, _, _),
  avg_latency = avg { api_response(req_id, _, latency, _) :
                      api_request(req_id, service_name, _, _, _),
                      latency },
  avg_latency > 2000.
\end{lstlisting}

\subsubsection{Mesh Compliance}

Universal quantification ensures all registered services are available
(line 211):

\begin{lstlisting}[language=Prolog,caption={Mesh compliance (rules.mg:211)}]
Decl mesh_compliant() :-
  forall(service_registry(S, _, _),
         service_status(S, /available)).
\end{lstlisting}

% ============================================================================
\section{Quality Controls}
\label{sec:quality}
% ============================================================================

Quality is enforced at multiple levels of the inference pipeline.

\subsection{Token Validation}

\subsubsection{EOS Checking Across Tokenizer Types}

The engine supports two tokenizer backends. EOS checking dispatches to the
correct one (\texttt{llama\_toon.zig} line 803):

\begin{lstlisting}[language=C,caption={Dual tokenizer EOS check (llama\_toon.zig:803)}]
const is_eos = if (self.gguf_tokenizer) |gt|
    gt.isEos(next_token)
else
    self.tokenizer.?.isEos(next_token);
\end{lstlisting}

EOS is checked for both primary sampling (line 804: \texttt{if (is\_eos) break})
and within accepted speculative drafts (lines 861--862).

\subsubsection{TOON Stop Sequences}

The engine also detects double-newline stop sequences for TOON-formatted output
(lines 906--912):

\begin{lstlisting}[language=C,caption={TOON stop detection (llama\_toon.zig:906--912)}]
if (n_out >= 2) {
    const last = output_tokens.items[n_out - 1];
    const prev = output_tokens.items[n_out - 2];
    const last_is_nl = (last == 13 or last == 10);
    const prev_is_nl = (prev == 13 or prev == 10);
    if (last_is_nl and prev_is_nl) break;
}
\end{lstlisting}

\subsection{Speculative Decode Acceptance Rates}

After all decode steps, the engine logs detailed speculation metrics
(\texttt{llama\_toon.zig} lines 927--933):

\begin{lstlisting}[language=C,caption={Speculation metrics (llama\_toon.zig:927--933)}]
if (spec_cycles > 0) {
    const alpha_pct: u64 = if (spec_drafted_total > 0)
        (spec_accepted_total * 100) / spec_drafted_total
    else 0;
    const eff_tps: u64 = if (decode_elapsed_ms > 0)
        (decoded_tokens * 1000) / decode_elapsed_ms
    else 0;
    std.log.info("spec_decode cycles={} drafted={} "
        ++ "accepted={} alpha={}% eff_tps={}", .{
        spec_cycles, spec_drafted_total,
        spec_accepted_total, alpha_pct, eff_tps,
    });
}
\end{lstlisting}

\subsection{Adaptive Disabling of Speculation}

When acceptance rates are too low (zero accepted tokens after the first cycle),
speculation is automatically disabled for the remainder of the request
(line 822). This prevents wasted batch-verification GPU cycles on inputs
where the draft model diverges from the target model.

\subsection{Timeout Guards}

Every decode iteration checks elapsed wall-clock time against a configurable
timeout (\texttt{llama\_toon.zig} lines 789--796):

\begin{lstlisting}[language=C,caption={Timeout guard (llama\_toon.zig:789--796)}]
const elapsed_ns = std.time.nanoTimestamp() - request_start_ns;
const elapsed_ms: u64 = @intCast(
    @max(elapsed_ns, 0) / std.time.ns_per_ms);
if (elapsed_ms >= request_timeout_ms) {
    std.log.err("generate timeout: elapsed_ms={} ...", .{...});
    return error.InferenceTimeout;
}
\end{lstlisting}

\subsection{DeltaNet State Management}

Memory safety for the recurrent DeltaNet state is enforced through:

\begin{enumerate}
  \item \textbf{Init with errdefer} (lines 363--366): If any allocation fails,
        previously allocated buffers are freed via \texttt{errdefer}:
\begin{lstlisting}[language=C,caption={errdefer safety (cuda\_forward.zig:363--366)}]
if (cuda.cuMemAlloc(&state.d_S, s_total) != .success)
    return error.OutOfMemory;
errdefer _ = cuda.cuMemFree(state.d_S);
\end{lstlisting}

  \item \textbf{Clear} (lines 410--415): Zero-fills all persistent state for
        new sequences via \texttt{cuMemsetD32}, preventing state leakage between
        requests.

  \item \textbf{Deinit} (lines 396--407): Iterates all device pointer fields,
        frees non-null pointers, and nulls them to prevent double-free.
\end{enumerate}

\subsection{Expert Cache with Hit/Miss Tracking}

The \texttt{ExpertCache} struct (\texttt{cuda\_forward.zig} lines 422--434)
maintains per-layer FP16 expert weight caches with explicit hit/miss counters.
On cache miss, Q4$\to$FP16 dequantization is performed and cached. On hit,
the expensive dequantization is skipped entirely, tracked via
\texttt{hits} and \texttt{misses} counters for operational monitoring.

\subsection{GPU Buffer Lifecycle}

All GPU buffer structs (\texttt{DartBatchBuffers} line 222,
\texttt{MoeBatchBuffers} line 260, \texttt{DeltaNetState} line 296) follow
a consistent pattern:
\begin{enumerate}
  \item Initialize pointers to zero.
  \item Allocate via \texttt{cuMemAlloc} with error propagation.
  \item In \texttt{deinit}, iterate all pointer fields, free non-null,
        set to zero.
\end{enumerate}

This pattern eliminates dangling pointers and provides deterministic cleanup
even under partial initialization failure.

\subsection{Mangle Quality Gates}

At the service mesh level, Mangle rules enforce:
\begin{itemize}
  \item \textbf{Routability}: Only services with \texttt{/available} status
        are routable (\texttt{rules.mg} line 10)
  \item \textbf{Latency thresholds}: Services with $>2000$ms average latency
        are flagged as slow (\texttt{rules.mg} line 97)
  \item \textbf{Mesh compliance}: Universal check that all registered services
        are available (\texttt{rules.mg} line 211)
  \item \textbf{Fallback routing}: Automatic reroute to \texttt{/chat}-capable
        services when primary is unavailable (\texttt{rules.mg} lines 32--38)
  \item \textbf{Intent fulfillment}: Explicit tracking of unfulfillable intents
        with reason strings (\texttt{rules.mg} lines 134--137)
\end{itemize}

% ============================================================================
\section{Conclusion}
\label{sec:conclusion}
% ============================================================================

This document has presented a private LLM inference engine built entirely in
Zig with hand-written CUDA kernels. The key architectural decisions are:

\begin{enumerate}
  \item \textbf{No Python dependency}: The entire forward pass, memory management,
        kernel dispatch, and speculative decoding are implemented in Zig 0.15.1.
  \item \textbf{DeltaNet hybrid support}: First-class support for Qwen3.5-style
        interleaved attention + linear attention architectures via 12 custom
        CUDA kernels compiled to PTX.
  \item \textbf{Multi-backend portability}: Unified abstraction across CUDA,
        Metal, WebGPU, and CPU backends with runtime capability negotiation.
  \item \textbf{Speculative decoding with adaptive control}: Dual-strategy draft
        generation (native MTP + n-gram fallback) with argmax acceptance and
        automatic disabling when speculation is unproductive.
  \item \textbf{Defense-in-depth quality controls}: Token validation, timeout
        guards, memory safety patterns, service health monitoring via Mangle
        rules, and operational metrics logging.
  \item \textbf{Kernel fusion}: Zero-CPU-touch pipeline from HTTP body to
        inference output, eliminating CPU-GPU round-trips for JSON parsing
        and tokenization.
\end{enumerate}

\end{document}
